\chapter{Inhalte laden – dafür brauchen Sie Internet}

Dies ist der Schritt, der aus einem leeren NOMAD ein Notfall-Wissensserver macht. Planen Sie Zeit ein: Das Laden dauert je nach Menge und Anschluss Stunden. Das System darf währenddessen nicht ausgeschaltet werden.

\kasten{nomadbraun}{Vorher prüfen}{%
\abhaken Der Rechner ist mit dem Internet verbunden (Kabel ist am sichersten).\par
\abhaken Das Netzteil steckt (nicht nur Akku).\par
\abhaken Freier Platz reicht für die geplante Menge (Tabelle in Kapitel~1). Bei der USB-SSD mit 512~GB sind alle deutschen Bereiche in „Standard“ (\SI{6,7}{GB}) und die deutsche Wikipedia „Beliebte Artikel“ (\SI{1,2}{GB}) bequem möglich.}

\section{Weg A: Der Schnellstart-Assistent (empfohlen)}
Klicken Sie auf der Startseite auf \menue{Schnellstart}. Der Assistent hat vier Schritte: \menue{Apps}, \menue{Karten}, \menue{Inhalte}, \menue{Überprüfen}.
\bild[0.85]{35-schnellstart-1}{Schnellstart, Schritt „Apps“: Wählen Sie die Funktionen, die Sie möchten. Für den Notfall-Wissensserver wählen Sie die \emph{Wissensbibliothek} (Wikipedia, Nachschlagewerke).}
\bild[0.85]{38-inhalte}{Schritt „Inhalte“: oben Wikipedia in mehreren Größen, darunter die Themenbereiche. Oben zeigt eine Leiste, wie viel Speicher Ihre Auswahl belegt.}
\bild[0.85]{39-medizin-stufen}{Jeder Themenbereich hat drei Stufen: \emph{Essential} (das Wichtigste, klein), \emph{Standard} und \emph{Comprehensive}. Die Größen stehen dabei.}
\bild[0.85]{41-ueberpruefen}{Schritt „Überprüfen“: Zusammenfassung mit der Gesamtgröße. Prüfen Sie sie, bevor Sie bestätigen.}

\merke{Beginnen Sie mit \textbf{Basis} in den drei deutschen Bereichen (zusammen etwa \SI{0,5}{GB}) und der deutschen Wikipedia „Schnellreferenz“ (\SI{0,14}{GB}). Das geht in wenigen Minuten und gibt Ihnen sofort ein arbeitsfähiges System. Danach, am besten über Nacht: „Standard“ und die deutsche Wikipedia „Beliebte Artikel“. Für Überleben und Vorsorge gibt es bisher nur englische Inhalte (Kategorie „Überleben \& Vorsorge (Englisch)“, etwa \SI{2,3}{GB} in der kleinsten Stufe).}

\section{Weg B: Einzelne Inhalte später nachladen}
Unter \menue{Einstellungen} $\rightarrow$ \menue{Content-Manager} finden Sie dieselben Kataloge. Dort laden Sie einzelne Inhalte nach oder löschen sie, wenn der Platz knapp wird.

\section{Wikipedia: Welche Größe?}
Wikipedia gibt es auf \textbf{Deutsch und Englisch}, jeweils in mehreren Größen. Die Größen sind die Anzeige in NOMAD. Es lässt sich genau eine Wikipedia auswählen.

\begin{center}\small
\renewcommand{\arraystretch}{1.25}
\begin{tabularx}{\linewidth}{@{}>{\RaggedRight\arraybackslash}p{5.6cm}rX@{}}
\toprule
\textbf{Auswahl} & \textbf{Größe} & \textbf{Geeignet für} \\
\midrule
Deutsche Wikipedia – Schnellreferenz & \SI{0,14}{GB} & die wichtigsten Artikel, kurz; schnelle Antworten \\
Deutsche Wikipedia – Beliebte Artikel & \SI{1,2}{GB} & wichtige Artikel in voller Länge, ohne Bilder \\
Deutsche Wikipedia – Komplett (Kompakt) & \SI{4,6}{GB} & alle Artikel, nur die Einleitung \\
Deutsche Wikipedia – Komplett (ohne Bilder) & \SI{18,1}{GB} & alle Artikel ohne Bilder \\
Deutsche Wikipedia – Komplett (Vollständig) & \SI{51,2}{GB} & alles mit Bildern; nur für große Platten \\
\midrule
Englische Wikipedia – Schnellreferenz bis Vollständig & \SIrange{0,3}{121}{GB} & dieselben Größenstufen, auf Englisch \\
\bottomrule
\end{tabularx}
\end{center}

Neben der Wikipedia enthalten die deutschen Kategorien unter anderem die Medizin-Enzyklopädie der Wikipedia (WikiMed), Wikibooks, Wikiversity, das deutsche Wörterbuch Wiktionary, das Kinderlexikon Klexikon, das Koch-Wiki, die iFixit-Reparaturanleitungen und die Bibliothek des Projekts Gutenberg.

\section{Karten}
Karten laden Sie unter \menue{Einstellungen} $\rightarrow$ \menue{Karten-Manager}. Dort gibt es zwei Wege:
\begin{description}[leftmargin=1.2em,style=nextline]
\item[Nach Land oder Region (empfohlen)] Knopf \menue{Länder auswählen}: Land ankreuzen, mit dem Regler die \emph{Zoomstufe} einstellen, den Knopf \menue{Download starten} drücken. NOMAD schneidet nur die gewählten Länder aus einer Weltkarte heraus und zeigt vorher die Größe an.
\item[Weltkarte] Knopf \menue{Weltkarte herunterladen}: der ganze Planet, etwa \SI{129}{GB}. Nur für große Platten.
\end{description}
Die Zoomstufe bestimmt den Detailgrad und damit die Größe. Beispiel Deutschland (Anzeige in NOMAD):

\begin{center}\small
\renewcommand{\arraystretch}{1.25}
\begin{tabular}{@{}rll@{}}
\toprule
\textbf{Zoomstufe} & \textbf{Detailgrad} & \textbf{Größe} \\
\midrule
10 & bis Stadtebene & \SI{0,12}{GB} (121~MB) \\
12 & Straßen größerer Orte & \SI{0,68}{GB} \\
15 & einzelne Straßen & \SI{5,6}{GB} \\
\bottomrule
\end{tabular}
\end{center}
\bild[0.85]{map01-karten-manager}{Der Karten-Manager: „Länder auswählen“ ist der empfohlene Weg.}
\bild[0.85]{map02-laender-waehlen}{Land ankreuzen, Zoomstufe wählen; unten steht die zu erwartende Größe.}
\bild[0.85]{map03-karte-deutschland}{Die Karte von Deutschland (Zoomstufe 10), offline angezeigt.}

\hinweis{Die Länderliste und die Beschriftung der Karte sind englisch (zum Beispiel „Germany“). Unter „Kuratierte Kartenregionen“ stehen nur US-Gebiete; für Deutschland nehmen Sie „Länder auswählen“.}

\section{Fortschritt beobachten}
\bild[0.85]{43-downloads}{Auf der Seite „Downloads“ sehen Sie, was gerade geladen wird und wie weit es ist. Nach einem Neustart geht das Laden an der Stelle weiter, an der es aufgehört hat.}

\section{Wann ist es fertig?}
Alles ist geladen, wenn die Downloads-Seite keinen Eintrag mehr zeigt und in der \menue{Wissensbibliothek} (Kiwix) alle gewählten Inhalte aufgelistet sind. Machen Sie danach \textbf{sofort den Offline-Test} (Kapitel~6), solange das Internet noch erreichbar ist.

\hinweis{Notieren Sie in Kapitel~9, welche Inhalte Sie geladen haben und mit welchem Datum. So wissen Sie später, wie aktuell Ihr Wissen ist.}
